\documentclass[11pt]{article}
\usepackage[margin=1in]{geometry}
\usepackage[T1]{fontenc}
\usepackage[utf8]{inputenc}
\usepackage{lmodern}
\usepackage{microtype}
\usepackage{amsmath,amssymb,amsthm,mathtools,bm}
\usepackage{booktabs,longtable,tabularx,array}
\usepackage{enumitem}
\usepackage{xcolor}
\usepackage{hyperref}
\hypersetup{
 colorlinks=true,
 linkcolor=blue!55!black,
 citecolor=green!45!black,
 urlcolor=blue!65!black,
 pdftitle={AP-E3: A Microscopic Level-Two-Magnitude Candidate},
 pdfauthor={Route F research note}
}
\newtheorem{theorem}{Theorem}[section]
\newtheorem{proposition}[theorem]{Proposition}
\newtheorem{lemma}[theorem]{Lemma}
\newtheorem{corollary}[theorem]{Corollary}
\theoremstyle{definition}
\newtheorem{definition}[theorem]{Definition}
\newtheorem{gate}[theorem]{Fail-closed gate}
\theoremstyle{remark}
\newtheorem{remark}[theorem]{Remark}
\newcommand{\CP}{\mathbb{CP}}
\newcommand{\CC}{\mathbb{C}}
\newcommand{\cO}{\mathcal O}
\newcommand{\cH}{\mathcal H}
\newcommand{\ii}{\mathrm i}
\newcommand{\dd}{\mathrm d}
\newcommand{\Tr}{\operatorname{Tr}}
\newcommand{\Span}{\operatorname{Span}}
\newcommand{\Sym}{\operatorname{Sym}}
\newcommand{\ket}[1]{\lvert #1\rangle}
\newcommand{\braket}[2]{\langle #1\vert #2\rangle}
\newcommand{\nopromote}{\texttt{physics\_promotion\_allowed=false}}

\title{AP-E3: A Microscopic Origin Candidate for Projective Level Magnitude Two\\
Two-Orbital Mott--Hund Locking, the Quadratic Veronese Line,\\
and the Exactly-Two UV Obstruction}
\author{Route F research note}
\date{14 July 2026}

\begin{document}
\maketitle

\begin{abstract}
This note asks which microscopic statement can produce the level magnitude
\(|k|=2\)
needed for a three-state projective carrier, rather than merely inserting that
integer in an effective action.  Large \(U(1)\) gauge invariance proves only
\(k\in\mathbb Z\) and cannot distinguish \(2\) from \(1\) or \(3\).  We then
construct a minimal conditional candidate.  Two spinful bosonic orbitals are
placed on a unit-occupancy Mott plateau and coupled ferromagnetically.  Their
low-energy Hilbert space is the symmetric spin-one triplet.  For a normalized
doublet \(s\), its oriented ground state is
\[
 \ket{s;2}=\ket{s}\otimes\ket{s},
 \qquad
 -\ii\braket{s;2}{\dd(s;2)}=2\bigl(-\ii s^\dagger\dd s\bigr).
\]
In standard algebraic-geometric conventions the ket eigenline is
\(L_{\rm ket}=\nu_2^*\cO_{\CP^2}(-1)=\cO_{\CP^1}(-2)\), whereas its dual
prequantum line is \(L_{\rm pre}=L_{\rm ket}^*=\cO_{\CP^1}(2)\).  Thus the
declared microscopic dimer derives \(|k|=2\) and
\(H^0(\CP^1,L_{\rm pre})\simeq\Sym^2\CC^2\) with dimension three, without a
bare level-two counterterm.  With the convention
\(A=-\ii\braket{s}{\dd s}\), however, the displayed microscopic kinetic term
gives signed level \(k=-2\); reversing the orientation gives \(k=+2\).

The construction is verified by a deterministic \(27/27\) audit.  At
\(U=4,\mu=1.5,J_H=1,h=0.2\), the sufficient interacting-plateau inequality
has margin \(2.025\), and the complete spinful occupancy problem has the
unique sector \((N_1,N_2)=(1,1)\), charge gap \(1.85\), Hund
spectrum \((-1/4,-1/4,-1/4,3/4)\), and a numerical Chern sequence converging
to \(2.000000501994128\).  The result is nevertheless conditional: the
repository does not yet derive why the UV field content contains exactly two
same-orientation constituents, why singleton or odd sectors are forbidden,
or why the resulting triplet is chiral Route-E family data.  The physical
orientation selecting the signed Route-E level is also not yet derived.
Consequently AP-E3 has a complete candidate-level magnitude derivation but
not a complete signed UV physics derivation, and \nopromote.
\end{abstract}
\tableofcontents

\section{Question, assumptions, and precise status}

AP-E1 derived the Hopf quotient and showed that a single projective doublet
naturally gives a ket line \(\cO(-1)\) and dual prequantum line \(\cO(1)\),
not their tensor squares.  AP-E2 froze the
independent discriminant--Killing theorem but did not produce a Berry level.
AP-E3 must therefore answer a logically separate question: what microscopic
Hilbert space, Hamiltonian, and gap produce the tensor-square Berry line, and
what additional principle makes that sector the lowest physical nonzero one?

The selected candidate uses the following declared data.
\begin{enumerate}[label=(A\arabic*)]
\item There are two orbitals \(r=1,2\), each with a spinor boson
  \(a_{rA}\), \(A=1,2\).
\item The full interacting core Hamiltonian lies on the unit-occupancy
  plateau.  A sufficient condition derived below is
  \(C:=\mu+h/2+J_H/4<U-J_H/8\).
\item A ferromagnetic Hund interaction locks the two spin-\(1/2\) moments into
  their symmetric triplet.
\item A weak orientation field selects a unique coherent ground-state line;
  all portal, drive, and thermal scales remain below the retained gaps.
\end{enumerate}
These are microscopic model inputs, not consequences of Route E.

\begin{theorem}[Conditional AP-E3 level-two-magnitude theorem]
Under assumptions (A1)--(A4), and while the many-body gap remains open, the
ground-state ket line and its dual prequantum line obey
\[
 L_{\rm ket}\simeq\cO_{\CP^1}(-2),\qquad
 L_{\rm pre}=L_{\rm ket}^*\simeq\cO_{\CP^1}(2),
\]
with \(c_1(L_{\rm ket})=-2\), \(c_1(L_{\rm pre})=2\), and
\[
 H^0(\CP^1,L_{\rm pre})\simeq\Sym^2\CC^2,\qquad \dim=3.
\]
Relative to \(S_k=\hbar k\int A\), the aligned Hamiltonian written below has
\(k=-2\); the oppositely oriented Hamiltonian has \(k=+2\).  Hence the model
derives \(|k|=2\), not the physical sign, without a bare level-two term.
\end{theorem}

\begin{gate}[What the theorem does not prove]
The theorem does not derive (A1), does not impose a fundamental superselection
against all single-constituent excitations, and does not identify a spin-one
internal multiplet with three chiral families.  A Mott gap makes singleton
sectors irrelevant to the strict low-energy EFT; it does not erase them from
the full UV Hilbert space.  A fundamental claim that magnitude two is the lowest
nonzero level therefore still requires an exactly-two UV rule or an
independent odd-sector prohibition.  It also does not derive which orientation
is the chiral Route-E convention, so the signed level remains a gate.
\end{gate}

\section{Large gauge invariance is an integrality theorem, not a selector}

Let \(s\in\CC^2\) satisfy \(s^\dagger s=1\), with redundancy
\(s\sim e^{\ii\alpha}s\).  The Hopf connection and curvature are
\begin{equation}
 A=-\ii s^\dagger\dd s,\qquad F=\dd A,\qquad
 \frac1{2\pi}\int_{\CP^1}F=1.
 \label{eq:hopf-data}
\end{equation}
For a closed worldline \(\gamma\), consider
\begin{equation}
 S_k=\hbar k\oint_\gamma A.
 \label{eq:level-action}
\end{equation}
Under \(s\mapsto e^{\ii\alpha}s\), one has \(A\mapsto A+\dd\alpha\), so a
large transformation of winding \(m\) changes the action by
\begin{equation}
 \Delta S_k=\hbar k\oint_\gamma\dd\alpha=2\pi\hbar k m.
\end{equation}
Single-valuedness of \(e^{\ii S/\hbar}\) gives exactly
\begin{equation}
 \boxed{k\in\mathbb Z.}
 \label{eq:integer-only}
\end{equation}
Levels \(1,2,3\) all pass this condition.  No argument based only on large
gauge invariance can select \(k=2\).

There are two common traps.
\begin{enumerate}
\item The transformation \(s\mapsto-s\) is already the gauge transformation
  \(\alpha=\pi\); by itself it does not require even \(k\).
\item Declaring a field to have gauge weight two need not double the faithful
  Hopf level.  If the Gauss law is \(2N=K\), then the physical occupation is
  \(N=K/2\); the minimal label \(K=2\) can still describe one Hopf quantum.
  A genuine \(\cO(2)\) line requires a tensor square, two additive occupied
  eigenlines, or another independently normalized topological term.
\end{enumerate}

An independent exact even-parity projector
\begin{equation}
 P_{\rm even}=\frac{1+(-1)^N}{2}
\end{equation}
would remove odd occupations.  For
\(E_N=UN(N-1)/2-\mu N\), the lowest nonzero even sector is uniquely \(N=2\)
when
\begin{equation}
 \frac U2<\mu<\frac{5U}{2}.
 \label{eq:even-window}
\end{equation}
But the parity projector and the plateau are new microscopic inputs;
Eq.~\eqref{eq:integer-only} does not generate them.

\section{The two-orbital Mott--Hund model}
\subsection{Microscopic action}

Let the number and spin operators be
\begin{equation}
 N_r=a_{rA}^\dagger a_{rA},\qquad
 \bm S_r=\frac12a_{rA}^\dagger\bm\sigma_{AB}a_{rB}.
 \label{eq:number-spin}
\end{equation}
Repeated spinor indices are summed.  The proposed microscopic action is
\begin{align}
 S_{\rm micro}
 &=\int\dd t\left[
  \ii\hbar\sum_{r,A}a_{rA}^\dagger\dot a_{rA}-H_{\rm micro}(\bm n)
 \right],\\
 H_{\rm micro}(\bm n)
 &=\sum_{r=1}^2\left[\frac U2N_r(N_r-1)-\mu N_r\right]
 -J_H\bm S_1\!\cdot\!\bm S_2
 -h\bm n\!\cdot\!(\bm S_1+\bm S_2)+H_{\rm portal}.
 \label{eq:micro-hamiltonian}
\end{align}
Here \(U,J_H,h>0\), and \(\bm n\in S^2\) is an adiabatic orientation
parameter.  The two-orbital structure is explicit rather than hidden in a
level coefficient.  Similar Mott/Hund and coherent-state techniques are
standard, but the specific exactly-two proposal here is a repository model,
not a theorem imported from the recent literature
\cite{ArovasAuerbach1988,OhZhang2023,ZhangEtAl2024Dimer}.

\subsection{Bare unit-occupancy plateau}

For one orbital the onsite energy is
\begin{equation}
 E_n^{(0)}=\frac U2n(n-1)-\mu n,\qquad n\in\mathbb Z_{\ge0}.
\end{equation}
If \(0<\mu<U\), then \(E_0^{(0)}-E_1^{(0)}=\mu>0\), and, for \(n\ge2\),
\begin{equation}
 E_n^{(0)}-E_1^{(0)}
 =(n-1)\left(\frac{Un}{2}-\mu\right)
 \ge(n-1)(U-\mu)>0.
\end{equation}
Thus \(n=1\) is the unique onsite ground occupation and the bare charge gap is
\begin{equation}
 \Delta_{\rm ch}^{(0)}=\min\{\mu,U-\mu\}.
 \label{eq:bare-charge-gap}
\end{equation}

\subsection{Exact interacting occupancy minimum and coercive tail}

Set \(H_{\rm portal}=0\) first.  For a single spinful bosonic mode with
occupation \(n_r\), Bose symmetry gives \(j_r=n_r/2\).  For \(J_H,h>0\), the
lowest state in the fixed \((n_1,n_2)\) sector has maximal total spin aligned
with \(\bm n\).  The exact core-Hamiltonian minimum is
\begin{align}
 E_{\min}(n_1,n_2)
 &=\sum_{r=1}^2\left[\frac U2n_r(n_r-1)-\mu n_r\right]
 -\frac{J_H}{4}n_1n_2-\frac h2(n_1+n_2).
 \label{eq:sector-minimum}
\end{align}

The bare window \(0<\mu<U\) is not by itself an interacting theorem.  For
example, \((U,\mu,J_H,h)=(1,0.99,3.9,0)\) satisfies both \(0<\mu<U\) and
\(4U>J_H\), but Eq.~\eqref{eq:sector-minimum} gives
\[
 E_{\min}(1,1)=-2.955,
 \qquad E_{\min}(2,2)=-5.86.
\]
The following inequality repairs that logical gap.

\begin{proposition}[Sufficient global interacting plateau]
Assume \(U,\mu,J_H,h>0\), set
\begin{equation}
 C:=\mu+\frac h2+\frac{J_H}{4},
 \label{eq:plateau-C}
\end{equation}
and suppose
\begin{equation}
 \boxed{C<U-\frac{J_H}{8}}
 \qquad\Longleftrightarrow\qquad
 \boxed{0<\mu<U-\frac h2-\frac{3J_H}{8}}.
 \label{eq:interacting-plateau-condition}
\end{equation}
Then \((n_1,n_2)=(1,1)\) is the unique global minimizer of
Eq.~\eqref{eq:sector-minimum} over \(\mathbb Z_{\ge0}^2\).  This condition is
sufficient, not claimed necessary.
\end{proposition}

\begin{proof}
Put \(x=n_1-1\) and \(y=n_2-1\), so \(x,y\in\mathbb Z_{\ge-1}\).  Direct
subtraction gives
\begin{equation}
 \begin{split}
 D(x,y)&:=E_{\min}(1+x,1+y)-E_{\min}(1,1)\\
 &=\frac U2\{x(x+1)+y(y+1)\}-C(x+y)-\frac{J_H}{4}xy.
 \end{split}
 \label{eq:plateau-difference}
\end{equation}
First take \(x,y\ge0\).  At fixed \(s=x+y\), Eq.~\eqref{eq:plateau-difference}
decreases as \(xy\) increases, so its minimum occurs at the most balanced
integer pair.  For \(s=2m\), \(m\ge1\), that minimum is
\begin{equation}
 D(m,m)=m\left[\left(U-\frac{J_H}{4}\right)m+U-2C\right]>0,
\end{equation}
because the bracket is increasing in \(m\) and at \(m=1\) equals
\(2(U-J_H/8-C)>0\).  Condition~\eqref{eq:interacting-plateau-condition}
also implies \(J_H<8U/3<4U\), hence \(U-J_H/4>0\).  For \(s=2m+1\), the
balanced value is
\begin{equation}
 D(m,m+1)=\left(U-\frac{J_H}{4}\right)(m^2+m)
 +\frac U2+\left(\frac U2-C\right)(2m+1).
\end{equation}
It starts at \(U-C>0\), and its successive difference is
\(2(m+1)(U-J_H/4)+U-2C>0\); hence it too is positive.

It remains to check the boundary allowed by \(x,y\ge-1\).  For \(y\ge0\),
\begin{equation}
 D(-1,y)=\frac U2y^2+\left(\frac U2-C+\frac{J_H}{4}\right)y+C.
\end{equation}
At \(y=0\) this is \(C>0\); at \(y=1\) it is \(U+J_H/4>0\), and for
\(y\ge1\) its successive difference
\(Uy+U-C+J_H/4\) is positive.  The case \(D(x,-1)\) is symmetric, while
\(D(-1,-1)=2C-J_H/4=2\mu+h+J_H/4>0\).  Every lattice point other than
\((x,y)=(0,0)\) therefore has strictly higher energy.
\end{proof}

Let \(N=n_1+n_2\).  From
\(n_1^2+n_2^2\ge N^2/2\) and \(n_1n_2\le N^2/4\),
\begin{equation}
 \boxed{
 E_{\min}(n_1,n_2)\ge
 \frac{4U-J_H}{16}N^2
 -\left(\frac U2+\mu+\frac h2\right)N.}
 \label{eq:coercive-bound}
\end{equation}
Hence \(4U>J_H\) makes the high-occupancy tail coercive.  Notice that
Eq.~\eqref{eq:interacting-plateau-condition} already implies this inequality.
A finite sector
enumeration together with Eq.~\eqref{eq:coercive-bound} proves the global
ground sector at any stated benchmark.

At
\begin{equation}
 (U,\mu,J_H,h)=(4,1.5,1,0.2)
 \label{eq:benchmark}
\end{equation}
one has
\[
 C=1.85<3.875=U-\frac{J_H}{8},\qquad
 (U-J_H/8)-C=2.025,
\]
and
\[
 E_{\min}(1,1)=-3.45,\qquad
 E_{\min}(1,0)=E_{\min}(0,1)=-1.60.
\]
The exact scan for \(N<8\) has the unique ground \((1,1)\) and interacting
charge gap \(1.85\).  For \(N\ge8\), the lower bound is increasing and is
already \(31.2\), excluding every omitted occupancy sector.

\section{Exact triplet projection and the orientation gap}

Inside \(N_1=N_2=1\), each orbital carries spin \(1/2\).  With
\(\bm S_{\rm tot}=\bm S_1+\bm S_2\),
\begin{equation}
 \bm S_1\cdot\bm S_2
 =\frac12\left(\bm S_{\rm tot}^2-\frac32\right).
 \label{eq:spin-addition}
\end{equation}
The Hund term has the exact spectrum
\begin{equation}
 E_{j=1}=-\frac{J_H}{4}\quad(3\text{-fold}),\qquad
 E_{j=0}=\frac{3J_H}{4}\quad(1\text{-fold}),
 \label{eq:hund-spectrum}
\end{equation}
with singlet gap \(\Delta_{\rm singlet}=J_H\).  The triplet projector is
\begin{equation}
 P_{j=1}=\frac{1+P_{12}}{2}=\frac{\bm S_{\rm tot}^2}{2},
 \qquad \Tr P_{j=1}=3.
 \label{eq:triplet-projector}
\end{equation}
Thus ferromagnetic locking is precisely projection onto \(\Sym^2\CC^2\).

For fixed \(\bm n\), the orientation term gives the four energies
\begin{equation}
 -\frac{J_H}{4}-h,\quad -\frac{J_H}{4},\quad
 -\frac{J_H}{4}+h,\quad \frac{3J_H}{4}.
 \label{eq:probe-spectrum}
\end{equation}
For \(h>0\), the ground state is unique and the triplet orientation gap is
\(h\).  If \(s(\bm n)\) is the normalized \(+1\) eigenvector of
\(\bm n\cdot\bm\sigma\), then
\begin{equation}
 \ket{\Omega(\bm n)}=\ket{s(\bm n)}\otimes\ket{s(\bm n)}.
 \label{eq:pair-ground}
\end{equation}
This is the microscopic line whose topology is computed next.

\section{Quadratic Veronese geometry and Berry factor two}

\subsection{The symmetric coherent state}

Write \(s=(z_0,z_1)^T\), with \(|z_0|^2+|z_1|^2=1\).  In the normalized
triplet basis,
\begin{equation}
 \ket{s;2}
 =z_0^2\ket{1,1}
  +\sqrt2z_0z_1\ket{1,0}
  +z_1^2\ket{1,-1}.
 \label{eq:veronese-state}
\end{equation}
Its norm is \((|z_0|^2+|z_1|^2)^2=1\), and it has no singlet component.
Projectively it is the quadratic Veronese embedding
\begin{equation}
 \nu_2:\CP^1\longrightarrow\CP^2,\qquad
 [z_0:z_1]\longmapsto[z_0^2:\sqrt2z_0z_1:z_1^2].
 \label{eq:veronese-map}
\end{equation}
Under \(s\mapsto e^{\ii\alpha}s\), the pair transforms as
\begin{equation}
 \ket{s;2}\mapsto e^{2\ii\alpha}\ket{s;2}.
 \label{eq:phase-two}
\end{equation}
This phase weight is a consequence of a tensor product, not a relabeling of
the gauge generator.

\subsection{Connection, curvature, and quantum metric}

Differentiating the tensor product gives
\begin{equation}
 \dd\ket{s;2}
 =\dd\ket{s}\otimes\ket{s}+\ket{s}\otimes\dd\ket{s}.
\end{equation}
Using \(\braket{s}{s}=1\),
\begin{equation}
 \braket{s;2}{\dd(s;2)}=2\braket{s}{\dd s},
\end{equation}
and therefore
\begin{equation}
 \boxed{A_2=2A_1,\qquad F_2=2F_1.}
 \label{eq:berry-addition}
\end{equation}
The Fubini--Study quadratic form is
\begin{equation}
 g(\dd\psi,\dd\psi)
 =\braket{\dd\psi}{\dd\psi}
  -|\braket{\psi}{\dd\psi}|^2.
\end{equation}
Substitution of the product derivative yields
\begin{equation}
 \boxed{g_2=2g_1.}
 \label{eq:metric-addition}
\end{equation}
Thus both symplectic and metric data are derived from the microscopic pair.
This is the standard coherent-state Berry construction applied to the tensor
product rather than assumed at the effective level
\cite{Berry1984,Perelomov1972APE3}.

\subsection{Patch proof, ket line, and dual prequantum line}

Use northern and southern unit sections
\begin{align}
 s_N&=\begin{pmatrix}\cos(\theta/2)\\e^{\ii\phi}\sin(\theta/2)\end{pmatrix},
 &A_N&=\frac{1-\cos\theta}{2}\dd\phi,\\
 s_S&=e^{-\ii\phi}s_N,
 &A_S&=-\frac{1+\cos\theta}{2}\dd\phi.
\end{align}
On the overlap, the pair transition is \(e^{-2\ii\phi}\).  Hence
\begin{equation}
 F_2=\sin\theta\,\dd\theta\wedge\dd\phi,
\end{equation}
The unitary connection on the ket line is
\(\nabla_{\rm ket}=\dd+\ii A_2\), whereas its dual has
\(\nabla_{\rm pre}=\dd-\ii A_2\).  Therefore
\begin{equation}
 \boxed{
 c_1(L_{\rm ket})=-\frac1{2\pi}\int_{\CP^1}F_2=-2,
 \qquad
 c_1(L_{\rm pre})=\frac1{2\pi}\int_{\CP^1}F_2=2.}
 \label{eq:chern-two}
\end{equation}
Equivalently, the transition function and Eq.~\eqref{eq:veronese-map} prove
\begin{equation}
 \boxed{
 L_{\rm ket}=\nu_2^*\cO_{\CP^2}(-1)=\cO_{\CP^1}(-2),
 \qquad
 L_{\rm pre}=L_{\rm ket}^*=\nu_2^*\cO_{\CP^2}(1)
 =\cO_{\CP^1}(2).}
 \label{eq:bundle-pullback}
\end{equation}

\begin{remark}[Why the dual is essential rather than notational]
The transition \(e^{-2\ii\phi}\) belongs to the algebraic ket eigenline
\(\cO(-2)\).  Holomorphic geometric quantization uses its positive dual
\(\cO(2)\).  This is also the sign distinction between the microscopic ket
kinetic term and the positive prequantum flux; it cannot be removed merely by
renaming the same line bundle.
\end{remark}

\subsection{Three-state quantization}

Holomorphic quantization of the dual prequantum line yields
\begin{equation}
 \cH_{|k|=2}=H^0(\CP^1,L_{\rm pre})=H^0(\CP^1,\cO(2))
 =\Span\{z_0^2,z_0z_1,z_1^2\}
 \simeq\Sym^2\CC^2,\qquad \dim\cH_{|k|=2}=3.
 \label{eq:borel-weil}
\end{equation}
The spin-one generators obey
\begin{equation}
 [J_x,J_y]=\ii J_z,\qquad J_x^2+J_y^2+J_z^2=2I_3.
\end{equation}
The same symmetric square therefore appears in the exact spectrum, Berry
line, and holomorphic state count.

\section{Low-energy action and robustness}

Adiabatic integration of gapped modes gives
\begin{equation}
 S_{\rm eff}[s]
 =-2\hbar\int A_1(s)-\int H_{\rm eff}(s)\,\dd t
 +O\left(\frac{\partial_t^2}{\Delta}\right),
 \label{eq:effective-action}
\end{equation}
where a conservative retained gap is
\begin{equation}
 \Delta=\min\{\Delta_{\rm ch},J_H,h,\Delta_{\rm extra}\}.
 \label{eq:retained-gap}
\end{equation}
Indeed \(\braket{s}{\dd s}=\ii A_1\), so the microscopic kinetic term gives
\(\ii\hbar\braket{s;2}{\dd(s;2)}=-2\hbar A_1\).  Thus, relative to
Eq.~\eqref{eq:level-action}, the aligned \(-h\bm n\cdot\bm S_{\rm tot}\)
Hamiltonian has \(k=-2\).  Replacing it by the anti-aligned orientation (for
example \(+h\bm n\cdot\bm S_{\rm tot}\)) reverses \(A_1\) and gives \(k=+2\).
The magnitude two is topological and cannot vary under a smooth perturbation
unless the line ceases to be isolated or its gap closes.  Which sign is the
physical Route-E chirality is not selected by the dimer itself.

At the benchmark, the verifier uses
\begin{equation}
 (\Lambda_{\rm portal},T,\hbar\omega)=(0.01,0.02,0.03),
\end{equation}
while the smallest retained gap is \(h=0.2\).  The ratio \(0.03/0.2=0.15\)
passes the declared \(0.2\) guard, whereas a portal scale \(0.22\) is rejected.
This is a benchmark separation check, not a universal error bound.

\section{Why the candidate still does not derive exactly two}

There are two distinct meanings of a level-two-magnitude origin.
\begin{enumerate}
\item \emph{Infrared origin:} after declaring two orbitals and working below
  all charge and singlet gaps, the retained ket/prequantum pair has Chern
  numbers \((-2,+2)\), hence \(|k|=2\).  This note proves that statement.
\item \emph{Fundamental selection:} the full UV theory explains why the
  number of same-orientation constituents is exactly two and prohibits every
  singleton or odd sector.  This note does not prove that statement.
\end{enumerate}

A single occupied orbital has ket/prequantum lines \(\cO(-1)\) and
\(\cO(1)\), respectively, and a two-state doublet.  The
Mott model makes charge-removal states heavy, so they are absent from the
strict IR EFT, but they remain states in the full microscopic theory.  To
claim that \(2\) is the fundamental lowest nonzero level, one needs at least
one additional mechanism:
\begin{enumerate}
\item an exact microscopic gauge constraint whose physical local operators
  are necessarily bilinears and whose odd sectors are absent;
\item a discrete gauge superselection with an independently derived even
  sector and a plateau selecting its first nonzero representation;
\item an anomaly or generalized-symmetry condition fixing a coefficient to
  two in a specified four-dimensional UV completion;
\item a lattice or unit-cell theorem deriving exactly two locked orbitals and
  forbidding isolated-orbital boundary states.
\end{enumerate}
This is the principal remaining AP-E3 blocker.

\section{Alternative microscopic branches}

\subsection{Filled-fermion determinant}

For \(r\) gapped fermion flavors, consider
\begin{equation}
 H_f(\bm n)=\sum_{a=1}^r\left[
 -\Delta_a\psi_a^\dagger(\bm n\cdot\bm\sigma)\psi_a
 +U_a(N_a-1)^2\right].
\end{equation}
If every flavor is singly occupied and remains gapped, the many-body ground
state is a Slater determinant of eigenlines.  Berry connections add:
\begin{equation}
 A_{\det}=\sum_a A_a,\qquad
 k_{\rm ind}=\sum_a C_a,\qquad C_a=\pm1.
 \label{eq:fermion-sum}
\end{equation}
Two same-orientation filled lines give \(|k|=2\).  One line gives \(|k|=1\),
and opposite orientations cancel to zero.  Flavor number, filling, and signs must
all be derived.  Recent experiments show that higher Chern sectors can be
engineered and tuned, supporting possibility rather than necessity
\cite{LeeEtAl2023}.

\subsection{Mixed WZW matching}

Let \(\omega_3\) and \(\omega_2\) be integrally normalized generators on the
two target factors.  A mixed Wess--Zumino--Witten term has schematic form
\begin{equation}
 S_{\rm WZW}=2\pi\hbar n\int_{Y_5}\omega_3(g)\wedge\omega_2(s),
 \qquad n\in\mathbb Z.
 \label{eq:mixed-wzw}
\end{equation}
For a localized configuration with baryon or skyrmion number \(B\), reduction
to its worldline gives
\begin{equation}
 S_{\rm wl}=\hbar nB\int_\gamma A,\qquad k=nB.
 \label{eq:wzw-reduction}
\end{equation}
Davighi and Lohitsiri construct a weakly coupled UV completion in which the
coefficient equals a color number and is tree-level exact
\cite{DavighiLohitsiri2024}.  Formally, \(n_c=2\) and \(B=1\) would give a
signed \(k=2\).  However, two-color theory has pseudoreal representations and a
different Pauli--Guersey global symmetry and coset
\cite{SuenagaEtAl2023}.  One cannot obtain a valid Route-E completion merely
by replacing \(n_c\) with two; the global quotient, anomaly, minimal soliton,
and mixed-WZW cohomology must be redone.

\subsection{Tangent and Dirac route}

The holomorphic tangent bundle satisfies
\begin{equation}
 T^{1,0}\CP^1\simeq\cO(2).
\end{equation}
If AP-E4 derives that a physical chiral fluctuation is tangent-valued, the
canonical Spin\({}^c\) or Dolbeault operator
\begin{equation}
 D_T=\sqrt2(\bar\partial_T+\bar\partial_T^\dagger)
\end{equation}
has index, with \(x=c_1(\cO(1))\),
\begin{equation}
 \operatorname{ind}D_T
 =\int_{\CP^1}\operatorname{ch}(\cO(2))\operatorname{Td}(T\CP^1)
 =\int_{\CP^1}(1+2x)(1+x)=3.
 \label{eq:dolbeault-index}
\end{equation}
This convention must be explicit.  An ordinary spin Dirac operator twisted
by \(\cO(p)\) instead has positive-chirality holomorphic sections
\(H^0(\cO(p-1))\), of dimension \(p\) for \(p\ge1\); at \(p=2\) that count is
two, not three.  Projective-space monopole spectra clarify these distinctions
but take the monopole level as input \cite{BykovSmilga2023APE3}.  The
tangent/Dirac branch belongs to AP-E4 and cannot be used backward to close
AP-E3.

\section{Failure modes and falsifiable gates}

\begin{longtable}{@{}>{\raggedright\arraybackslash}p{0.25\textwidth}
                       >{\raggedright\arraybackslash}p{0.66\textwidth}@{}}
\toprule
Failure & Consequence and required check\\
\midrule
\endfirsthead
\toprule
Failure & Consequence and required check\\
\midrule
\endhead
\(\mu\notin(0,U)\) & The bare onsite ground changes to \(N=0\) or
\(N\ge2\); repeat the complete occupancy minimization.\\
Interacting plateau inequality fails & The bare condition \(0<\mu<U\) no
longer proves \((1,1)\); minimize Eq.~\eqref{eq:sector-minimum} globally.\\
\(J_H<0\) & The singlet is the unique ground state and carries \(k=0\).\\
No locking & The target is \(\CP^1\times\CP^1\), not its diagonal
\(\CP^1\).\\
Opposite Berry orientations & The two Chern numbers cancel, yielding
\(k=0\).\\
Singleton allowed & A physical \(|k|=1\) doublet remains; the fundamental
exactly-two claim fails.\\
Orientation not derived & The declared aligned action has \(k=-2\); obtaining
the desired \(k=+2\) requires a physical orientation/chirality argument.\\
\(h=0\) & The triplet ground space is degenerate, so there is no unique
Abelian ground-state line for the displayed Berry connection.\\
Gap-scale collision & If portal, temperature, or drive scales reach a charge,
singlet, orientation, or extra-state gap, the one-line adiabatic EFT is not
controlled.\\
Internal triplet only & A three-state spin multiplet is not automatically
three chiral families; construct the four-dimensional Dirac operator and
portal.\\
\bottomrule
\end{longtable}

The strongest next AP-E3 tests are:
\begin{enumerate}
\item construct an explicit exactly-two gauge constraint or unit-cell
  symmetry and prove the absence of boundary singleton sectors;
\item embed Eq.~\eqref{eq:micro-hamiltonian} in an anomaly-consistent
  four-dimensional model;
\item compute the many-body Chern number with a lattice-link or Wilson-loop
  algorithm under every portal deformation up to the first gap closing
  \cite{FukuiHatsugaiSuzuki2005};
\item prove that the Route-E coupling maps the triplet to chiral family data
  and does not introduce opposite-chirality partners.
\end{enumerate}

\section{The 27-check reproducible audit}

The executable verifier is
\begin{center}
\path{route_f/code/verify_ap_e3_level_two_microscopic.py}.
\end{center}
From the repository root, run
\begin{center}
\texttt{python3 route\_f/code/verify\_ap\_e3\_level\_two\_microscopic.py}.
\end{center}
It writes JSON and Markdown reports in \path{route_f/output/}.  The seed is
\(20260714\).  Its checks are grouped as follows.

\begin{center}
\begin{tabularx}{\textwidth}{@{}lclX@{}}
\toprule
Group & Count & Status & Content\\
\midrule
M0 provenance & 2 & pass & Input hashes and green, non-promoting AP-E2 gate.\\
M1 Mott & 4 & pass & Bare window, false-inference counterexample, proved
interacting plateau inequality, full minimum, and coercive tail.\\
M2 Hund & 7 & pass & Hermiticity, \(SU(2)\), triplet/singlet spectrum,
projector, gaps, sign reversal, and 64 orientations.\\
M3 Veronese & 6 & pass & Normalization, phase weight, Berry and metric factor
two, ket/dual transition signs, and prequantum Chern quadrature.\\
M4 representation & 2 & pass & Spin-one Casimir and large-gauge
non-selection.\\
M5 failures & 3 & pass & Opposite orientation, singleton, and unlocked target.\\
M6 EFT & 2 & pass & Positive hierarchy and rejected bad portal.\\
M7 boundary & 1 & pass & Magnitude proof complete; signed orientation and UV
physics closure false.\\
\midrule
Total & 27 & pass & \\
\bottomrule
\end{tabularx}
\end{center}

For \(J_H=1\), exact diagonalization gives
\begin{equation}
 \operatorname{spec}H_{\rm Hund}=(-0.25,-0.25,-0.25,0.75),
\end{equation}
while the antiferromagnetic sign gives
\((-0.75,0.25,0.25,0.25)\).  Over 64 seeded orientations at \(h=0.2\), the
maximum spectral residual is \(8.88\times10^{-16}\).  Over 256 random
normalized sections, the maximum Veronese or norm residual is
\(9.99\times10^{-16}\), the Berry factor-two residual is
\(2.24\times10^{-16}\), and the metric residual is
\(2.22\times10^{-16}\).

For \(N_\theta\) midpoint cells, direct curvature quadrature gives
\begin{equation}
 c_1^{(N_\theta)}
 =\frac{\pi}{N_\theta}
  \sum_{j=0}^{N_\theta-1}
  \sin\left(\frac{(j+1/2)\pi}{N_\theta}\right).
\end{equation}
The deterministic sequence is
\begin{center}
\begin{tabular}{@{}rr@{}}
\toprule
\(N_\theta\) & \(c_1^{(N_\theta)}\)\\
\midrule
20 & 2.0020576483\\
80 & 2.0001285163\\
320 & 2.0000080319\\
1280 & 2.0000005020\\
\bottomrule
\end{tabular}
\end{center}
Its monotone second-order convergence to the dual-prequantum value two verifies
Eq.~\eqref{eq:chern-two} numerically; it does not replace the patch proof.

\section{Recent-literature boundary}

Recent primary work supplies useful ingredients but not the missing theorem:
\begin{itemize}
\item projective monopole spectra and twisted Dolbeault complexes support the
  AP-E4 spectral audit, while treating the level as an input
  \cite{BykovSmilga2023APE3};
\item tunable higher Chern numbers demonstrate realizability but also show
  that a Hamiltonian can select different integers in different regimes
  \cite{LeeEtAl2023};
\item mixed-WZW generalized-symmetry matching can protect an integer
  coefficient, but its global symmetry and color content must match the UV
  completion \cite{DavighiLohitsiri2024};
\item coupled-dimer coherent-state methods make singlet/triplet structure
  calculable, but antiferromagnetic pairing itself does not imply the
  ferromagnetic diagonal line or \(|k|=2\) \cite{ZhangEtAl2024Dimer}.
\end{itemize}
No cited paper proves the repository-specific chain
\begin{equation}
 \begin{gathered}
 \text{exactly two constituents}\Longrightarrow\text{no singleton},\\
 \Longrightarrow\{\cO(-2)_{\rm ket},\cO(2)_{\rm pre}\}
 \Longrightarrow \text{three chiral Route-E families}.
 \end{gathered}
\end{equation}

\section{Conclusion and next decisions}

Within the declared two-orbital Mott--Hund model, AP-E3 now has a complete
first-principles magnitude chain:
\begin{equation}
 \begin{gathered}
 C=\mu+h/2+J_H/4<U-J_H/8
 \ \Longrightarrow\ (N_1,N_2)=(1,1),\\
 J_H>0\ \Longrightarrow\ \Sym^2\CC^2\text{ triplet},\\
 \ket{s}\otimes\ket{s}\ \Longrightarrow\ A_2=2A_1,\\
 L_{\rm ket}=\nu_2^*\cO(-1)=\cO(-2),\qquad
 L_{\rm pre}=L_{\rm ket}^*=\cO(2),\\
 S_{\rm aligned}=-2\hbar\int A_1,\qquad
 S_{\rm reversed}=+2\hbar\int A_1,\\
 |k|=2\ \Longrightarrow\quad
 \dim H^0(\CP^1,\cO(2))=3.
 \end{gathered}
 \label{eq:final-chain}
\end{equation}
The analytic occupancy theorem and \(27/27\) numerical audit make this more than
a formal analogy.  The construction also has sharp counterexamples:
antiferromagnetic exchange gives a \(k=0\) singlet, opposite orientations
cancel, no locking leaves a product target, and a singleton restores
\(|k|=1\).  The aligned sign \(k=-2\) also makes the Route-E chirality choice
an explicit unresolved input rather than a hidden convention.

The candidate-level task is therefore complete.  The full AP-E3 physics task
is not.  Three viable research choices remain:
\begin{enumerate}
\item \textbf{Preferred minimal route:} derive a microscopic gauge or
  unit-cell rule that makes the two-orbital composite fundamental and removes
  odd sectors, then construct the Route-E portal.
\item \textbf{Topological UV route:} build the correct \(n_c=2\) mixed-WZW
  completion and redo its global-symmetry and anomaly analysis.
\item \textbf{Independent AP-E4 route:} derive a tangent-valued chiral Dirac
  mode and its complete partner spectrum, without using the desired count as
  an input.
\end{enumerate}
Until one of these closes its own gates,
\begin{equation}
 \boxed{\begin{gathered}
 \texttt{candidate\_level\_magnitude\_derivation\_complete=true},\\
 \texttt{signed\_route\_e\_level\_selected=false},\\
 \texttt{ap\_e3\_physics\_closed=false},\\
 \texttt{physics\_promotion\_allowed=false}.
 \end{gathered}}
\end{equation}

\bibliographystyle{unsrt}
\bibliography{route_f/tex/ap_e3_level_two_microscopic}

\end{document}
